% English translation — Grothendieck fonds, folder 115, batch 1 (leaves 1-14).
% Derived from batch-01.fr.tex, never from the manuscript: two independent
% readings of the same hand would diverge, and nothing would say which was
% right. The apparatus carries over unchanged, in the same places.
% Pass: Fable 5 (claude-fable-5), 2026-08-08 — first pass, unchecked against the leaves by a human.
\documentclass[11pt,a4paper]{article}
% Le préambule commun à toutes les transcriptions du fonds Grothendieck.
%
% Il tient en une idée : **l'incertitude se compose, elle ne se résout pas.**
% Une transcription qui lisse un mot douteux a détruit la seule chose pour
% laquelle on la faisait — un lecteur doit pouvoir savoir, à chaque endroit,
% ce qui était sur le feuillet et ce qui a été ajouté. Les six macros qui
% suivent sont donc l'appareil critique, et non de la décoration.
%
% Les mêmes commandes sont reconnues par `scripts/render.mjs`, qui produit la
% vue de lecture du site. Une macro ajoutée ici doit l'être aussi là-bas,
% faute de quoi le rendu échouera bruyamment — ce qui est voulu : un rendu
% silencieusement faux est pire qu'un rendu absent.

\ProvidesPackage{grothendieck}[2026/08/08 Transcriptions du fonds Grothendieck]

\usepackage{amsmath,amssymb,amsthm}
\usepackage{mathrsfs}
\usepackage{stmaryrd}
% Les diagrammes commutatifs sont partout dans ce fonds : c'est la langue
% même de la géométrie algébrique de Grothendieck, et une transcription qui
% les rendrait en prose ne transcrirait rien.
\usepackage{tikz-cd}
% Français par défaut — c'est la langue des feuillets — mais l'anglais est
% chargé aussi : la traduction et le résumé sont des documents anglais, et la
% typographie française (espaces avant les deux-points) y serait une faute.
% Ces documents basculent par \selectlanguage{english} en tête de corps.
\usepackage[english,french]{babel}
\usepackage{fontspec}
\usepackage{geometry}
\geometry{a4paper,margin=25mm,marginparwidth=28mm}
\usepackage{marginnote}
\usepackage{xcolor}
% ulem, not soul. `soul`'s \st and \ul are built on a character-by-character
% scan that XeLaTeX's font machinery breaks: the document fails with an
% undefined control sequence rather than degrading. ulem does the same two jobs
% and is Unicode-engine safe. `normalem` stops it hijacking \emph, which the
% transcriptions use for Grothendieck's own emphasis.
\usepackage[normalem]{ulem}
\usepackage{hyperref}
\hypersetup{colorlinks=true,linkcolor=black,urlcolor=black,pdfborder={0 0 0}}

% --- Métadonnées du lot -----------------------------------------------------
% Reprises telles quelles par le script de rendu, qui les affiche en tête de
% la vue de lecture. Elles fixent ce que le fichier prétend couvrir : sans
% elles, une transcription flotte sans point d'attache dans un fonds de seize
% mille feuillets.

\newcommand{\folder}[1]{\gdef\grothFolder{#1}}
\newcommand{\batch}[1]{\gdef\grothBatch{#1}}
\newcommand{\leaves}[2]{\gdef\grothFirst{#1}\gdef\grothLast{#2}}
\newcommand{\foldertitle}[1]{\gdef\grothFoldertitle{#1}}
\gdef\grothFolder{?}\gdef\grothBatch{?}\gdef\grothFirst{?}\gdef\grothLast{?}\gdef\grothFoldertitle{}

% --- Le résumé --------------------------------------------------------------
%
% Il ouvre la lecture modernisée, sous le titre « Résumé », et il est composé
% en petit. Ce n'est pas un ornement : c'est le seul passage du document qui
% ne soit pas des mathématiques, et un lecteur doit pouvoir le reconnaître
% avant de l'avoir lu — puis le sauter s'il connaît déjà le sujet, ou s'y
% arrêter s'il ne le connaît pas. Le filet qui le suit marque la reprise du
% corps.

\newenvironment{resume}
  {\small\setlength{\parskip}{0.45em}}
  {\par\medskip\hrule\medskip}

% --- Mots-clés ---------------------------------------------------------------
%
% En anglais, en clôture du résumé de la lecture modernisée : le vocabulaire
% moderne sous lequel chercher ce que les feuillets construisent. Le manifeste
% du site (`npm run manifest`) les extrait pour étiqueter la cote entière —
% cette ligne est la source unique des tags, il n'y a pas de fichier à côté.
\newcommand{\keywords}[1]{\par\smallskip\noindent
  {\footnotesize\textsc{Keywords} --- \textit{#1}}\par}

% --- L'appareil critique ----------------------------------------------------

% Le numéro de feuillet, en marge. C'est l'ancre qui permet de régler un
% désaccord sur une formule en regardant *un* feuillet plutôt qu'une chemise
% de six cents. Le site s'en sert aussi pour tourner les pages du fac-similé
% quand on fait défiler la transcription.
\newcommand{\leaf}[1]{%
  \marginnote{\footnotesize\color{gray!70}\textsf{#1}}%
  \label{leaf:#1}\ignorespaces}

% Illisible : on ne devine pas. Un mot inventé qui se lit comme les autres est
% le pire résultat possible.
\newcommand{\ill}{\textcolor{red!60!black}{[\,\ldots\,]}}

% Lecture douteuse : le mot est donné, et signalé comme incertain.
\newcommand{\uncertain}[1]{\uline{#1}}

% Ajout de l'éditeur — une lettre manquante, un mot sous-entendu.
\newcommand{\add}[1]{\textcolor{blue!50!black}{[#1]}}

% Biffé par Grothendieck lui-même. Ce qu'il a rayé fait partie du document :
% une rature dit souvent où la pensée a changé de direction.
\newcommand{\struck}[1]{\sout{#1}}

% Note du transcripteur, sur l'état matériel du feuillet.
\newcommand{\note}[1]{\par\smallskip{\small\color{gray!60!black}\textit{[#1]}}\par\smallskip}

% Note marginale de Grothendieck, distincte du corps du texte.
\newcommand{\marginal}[1]{\par\smallskip\noindent\hspace{1em}%
  {\small\color{gray!60!black}#1}\par\smallskip}

% --- Titre ------------------------------------------------------------------

\AtBeginDocument{%
  \begin{center}
    {\large\bfseries Fonds Alexandre Grothendieck --- cote n\textsuperscript{o}~\grothFolder}\\[2pt]
    {\itshape\grothFoldertitle}\\[4pt]
    {\small feuillets \grothFirst--\grothLast\ (lot \grothBatch)}\\[2pt]
    {\footnotesize\color{gray!60!black}Transcription établie d'après le fac-similé numérisé
     par l'Université de Montpellier.\\
     Les lectures incertaines sont soulignées, les illisibles notées [\,\ldots\,],
     les ajouts entre crochets.}
  \end{center}
  \vspace{1em}}

\endinput


\folder{115}
\batch{1}
\leaves{1}{14}
\foldertitle{« Correspondances » fonctorielles. Dualité des topos — English translation}

\begin{document}
\selectlanguage{english}

\begin{summary}
These leaves — several written across the backs of IBM job listings dated
June 1982, which is how the archivists date the folder — pursue one question:
how much of a category is captured by the two kinds of shadow its objects
cast?

Fix a small category $A$ and think of its objects as shapes. Each object can
be probed from the left, by the maps arriving into it, or from the right, by
the maps leaving it. Recording all arrivals turns an object into a
\emph{presheaf}; Grothendieck writes $\widehat{A}$ for the category of these.
Recording departures gives a \emph{copresheaf}, living in $A^{\vee}$. These
are the two shadow-worlds, and both are topoi — hence the title.

The first leaves set up the ``functorial correspondences'': a single object
$X$ over a product $A \times B$ acts like an integral kernel $k(x,y)$ in
analysis, transforming shadows on $A$ into shadows on $B$ and back — the
transforms ${}^{s}X$ and ${}^{d}X$ of leaf 3, with the pairing $X(L,M)$.
Leaf 5 arranges eight functor categories in a wheel, all equivalent: eight
ways of reading the same kernel, one per choice of variance.

The duality itself is on leaf 7. A pairing $\Phi \star_{A} \Psi$ between
presheaves and copresheaves gives an adjunction between $\widehat{A}$ and
$A^{\vee}$ — each is ``functions on'' the other, the way a vector space $V$
and its dual $V^{*}$ are in linear algebra. His \emph{Question} — when are
the adjunction maps isomorphisms? — is the reflexivity question of linear
algebra, one level up. Today this is called Isbell duality. On leaves 10--11
he assembles triples $(H_{*}, H^{*}, \alpha)$ — a left shadow, a right
shadow, and a composition law binding them — a construction now known as the
Isbell envelope; and on leaf 13 he relates the self-dual core to
$\mathrm{Kar}(A)$, the category obtained by splitting idempotents, in modern
terms the Cauchy completion.

Nothing here was published. The notes sit among the material around
\emph{Pursuing Stacks} (1983), and the circle of ideas is now standard: the
reader who wants today's language should look up \emph{profunctor},
\emph{Isbell duality} and \emph{Cauchy completion}.
\end{summary}

\section*{Functorial correspondences (leaves 1 to 5)}

\leaf{1}
Title inked on the folder: « Correspondances » fonctorielles — Dualité des
topos. \note{folder cover}

\leaf{2}
Adjunction formulas, at the head of the leaf:
\[
H_{A}\bigl(f^{*}(\Psi_{\varepsilon}),\ \Phi^{\varepsilon}\bigr)
\;\simeq\;
H_{B}\bigl(\Psi_{\varepsilon},\ f^{\varepsilon}_{!}(\Phi^{\varepsilon})\bigr)
\]
\note{uncertain reading of the $\varepsilon$ decorations and of the placement
of indices}
\[
H_{A}\bigl(f^{*}(b),\ a\bigr) \;\simeq\; H_{B}\bigl(b,\ f^{\varepsilon}(a)\bigr)
= \mathrm{Hom}_{B}\bigl(b,\ f(a)\bigr),
\qquad
f^{\varepsilon}(a) = \varinjlim \cdots
\]
\note{the index system of the colimit is illegible. The rest of the leaf —
a circle of functor categories and small diagrams — is a struck-out first
state of the wheel of leaf 5; the readings are too uncertain to transcribe}

\leaf{3}
\note{author's page 1}
At the head, right: $\underline{\mathrm{Hom}}_{!}$: full subcategory of
$\underline{\mathrm{Hom}}$ formed by the functors commuting with
\uncertain{inductive limits}; $\underline{\mathrm{Hom}}^{!}$: full
subcategory of $\underline{\mathrm{Hom}}$ formed by the functors commuting
with \uncertain{projective limits}.

Left margin: case $B = A^{\circ}$: whence the canonical element
$H = H_{A} \in (A \times A^{\circ})^{\wedge}$; \uncertain{the canonical
extensions} $\widehat{A}^{\circ} \hookrightarrow A^{\vee}$, \ldots
\note{the rest of the marginal column is too heavily crossed out to
transcribe}

The central diagram, simplified — the original is a trellis of a dozen
arrows marked $\approx$, several of them struck out:

\begin{tikzcd}[column sep=small, row sep=small, nodes={font=\scriptsize}]
\underline{\mathrm{Hom}}(B^{\circ}, \widehat{A})
  & (A\times B)^{\wedge} \arrow[l, "\approx"'] \arrow[r, "\approx"] \arrow[d, "\approx"]
  & \underline{\mathrm{Hom}}(A^{\circ}, \widehat{B}) \arrow[d, "\approx"] \\
  & \underline{\mathrm{Hom}}^{!!}\bigl(\widehat{A}^{\circ} \times \widehat{B}^{\circ},\ (\mathrm{Ens})\bigr) = \mathrm{Bifais}(\widehat{A}, \widehat{B})
  & \underline{\mathrm{Hom}}^{!}(\widehat{A}^{\circ}, \widehat{B}) \\
  & \underline{\mathrm{Hom}}_{!}(A^{\vee}, \widehat{B}) \arrow[u, "\approx"']
  &
\end{tikzcd}

The systematic set-up: $A \times B$, with projections
$p_{B} : A\times B \to A$ and $p_{A} : A\times B \to B$.
\note{the projection's index names the factor being removed — the only
reading under which the following formulas agree with one another}
\[
\underline{\mathrm{Hom}}_{A}(X,Y) \overset{\mathrm{def}}{=}
   p_{B*}\,\underline{\mathrm{Hom}}_{A\times B}(X,Y),
\qquad
\underline{\mathrm{Hom}}_{B}(X,Y) \overset{\mathrm{def}}{=}
   p_{A*}\,\underline{\mathrm{Hom}}_{A\times B}(X,Y)
\]
\[
\mathrm{Hom}_{A\times B}(X,Y)
\simeq \Gamma_{A\times B}\,\underline{\mathrm{Hom}}_{A\times B}(X,Y)
\simeq \Gamma_{A}\,\underline{\mathrm{Hom}}_{A}(X,Y)
\simeq \Gamma_{B}\,\underline{\mathrm{Hom}}_{B}(X,Y),
\qquad X, Y \in (A\times B)^{\wedge}
\]
\[
{}^{s}X(L) = \underline{\mathrm{Hom}}_{B}\bigl(p_{B}^{*}L,\ X\bigr)
           = p_{A*}\,\underline{\mathrm{Hom}}\bigl(p_{B}^{*}L,\ X\bigr),
\qquad
{}^{d}X(M) = \underline{\mathrm{Hom}}_{A}\bigl(p_{A}^{*}M,\ X\bigr)
           = p_{B*}\,\underline{\mathrm{Hom}}\bigl(p_{A}^{*}M,\ X\bigr)
\]
(in the margin: $L \in \widehat{A}$, $M \in \widehat{B}$)
\[
X(L,M) = \mathrm{Hom}\bigl(L \boxtimes M,\ X\bigr)
\simeq \mathrm{Hom}\bigl(L,\ {}^{d}X(M)\bigr)
\simeq \mathrm{Hom}\bigl(M,\ {}^{s}X(L)\bigr),
\qquad\text{where } L \boxtimes M = p_{B}^{*}L \times p_{A}^{*}M .
\]

``These formulas are particularly interesting for
$L = a \in \mathrm{Ob}\,A$, $M = b \in \mathrm{Ob}\,B$.''
\note{uncertain reading of the two right-hand members}

``Let $X$; let us define $B^{\circ} \to \widehat{A}$ (namely
$b \mapsto {}^{d}X(b) = \underline{\mathrm{Hom}}_{A}(p_{A}^{*}(b), X)$)
\ldots{} one finds $B^{\circ} \to \mathrm{Ens}$: it is \ill{}, which can
also be written $\underline{\mathrm{Hom}}_{B}(p_{B}^{*}(L), X) = {}^{s}X(L)$.
\uncertain{Dual} interpretation for ${}^{d}X(M)$.''
\note{interlined passage, reading partly conjectural}

``$(L,M) \mapsto L \boxtimes M : \widehat{A} \times \widehat{B} \to
(A\times B)^{\wedge}$ is fully faithful.''

``the \ill{} which is the canonical functor \ldots{} said to be ``of
Yoneda'' for $\widehat{A} \times \widehat{B}$''
\[
\underline{\mathrm{Hom}}_{A\times B}\bigl(L\boxtimes M',\ L'\boxtimes N'\bigr)
\simeq \underline{\mathrm{Hom}}_{A}(L,L') \boxtimes
       \underline{\mathrm{Hom}}_{B}(M',N')
\qquad \text{e.g.}
\]
\[
\underline{\mathrm{Hom}}_{A\times B}\bigl(L\boxtimes M',\ L'\boxtimes e_{B}\bigr)
\simeq \underline{\mathrm{Hom}}_{A}(L,L') \boxtimes e_{B}
= p_{B}^{*}\,\underline{\mathrm{Hom}}_{A}(L,L')
\]
in particular, taking $L' = e_{A}$: \ldots
\note{the rest of the line is illegible}

$\underline{\mathrm{Hom}}_{A\times B}\bigl(p_{A}^{*}(M), \cdots\bigr)
\Leftarrow p_{B}^{*}\cdots$ (the arrows being the canonical ones)
\note{line partly illegible}

Setting $X^{M} \overset{\mathrm{def}}{=}
\underline{\mathrm{Hom}}_{A\times B}\bigl(p_{A}^{*}M,\ X\bigr)$, one thus
finds $(p_{B}^{*}L)^{M} \simeq p_{B}^{*}L\cdots$
\note{end of formula uncertain}

\leaf{5}
\note{author's page 2}
Top margin: $A^{\circ}, \widehat{A}^{\circ}, A^{\vee}$;
$B^{\circ}, \widehat{B}^{\circ}, B^{\vee}$;
$\bigl|\widehat{A}^{\circ} \times A^{\vee}\bigr| \times
 \bigl|\widehat{B}^{\circ} \times B^{\vee}\bigr| \to \mathrm{Ens}$
\note{decorations ``(4 co)'', ``(6 (co))'' not elucidated}

The wheel of eight functor categories around
$(A\times B)^{\wedge} = \underline{\mathrm{Hom}}(A^{\circ}, B^{\circ};\,
\mathrm{Ens})$:

\begin{tikzcd}[column sep=tiny, row sep=small, nodes={font=\scriptsize}]
& \underline{\mathrm{Hom}}^{!!}(\widehat{A}^{\circ}\times\widehat{B}^{\circ};\ \mathrm{Ens}) \arrow[dl, "\approx"'] \arrow[dr, "\approx"] & \\
\underline{\mathrm{Hom}}^{!}(\widehat{B}^{\circ},\widehat{A}) \arrow[d, "\approx"'] & & \underline{\mathrm{Hom}}^{!}(\widehat{A}^{\circ},\widehat{B}) \arrow[d, "\approx"] \\
\underline{\mathrm{Hom}}(B^{\circ},\widehat{A}) \arrow[d, "\approx"'] & (A\times B)^{\wedge} & \underline{\mathrm{Hom}}(A^{\circ},\widehat{B}) \arrow[d, "\approx"] \\
\underline{\mathrm{Hom}}_{!}(B^{\vee},\widehat{A}) \arrow[dr, "\approx"'] & & \underline{\mathrm{Hom}}_{!}(A^{\vee},\widehat{B}) \arrow[dl, "\approx"] \\
& \underline{\mathrm{Hom}}_{!!}(A^{\vee}\times B^{\vee};\ \mathrm{Ens}) &
\end{tikzcd}

``all the arrows of the diagram \ldots{} are equivalences of categories''
\note{the quadrants of the original are marked $\wedge\wedge$, $\wedge\vee$,
$\vee\vee$, $\vee\wedge$, and the outer arcs $(\widehat{A}\to)$,
$(\widehat{B}\to)$, $(A^{\vee}\to)$, $(B^{\vee}\to)$; the two left-hand
nodes, partly hidden by a tear, are of uncertain reading}

At the bottom — index families: ``$I$ an index set, $(A_i)_{i\in I}$ a
family of (small) categories. If $I' \subset I$, set
$A_{I'} = \prod_{i\in I'} A_i$.''
\[
I = \coprod_{j\in J} I_{j}, \qquad J = J^{+} \amalg J^{-},
\qquad I^{+} = \coprod_{j\in J^{+}} I_{j}, \quad
I^{-} = \coprod_{j\in J^{-}} I_{j}
\]
``\uncertain{I assume} $I$, $J$ finite to be « plus safe ».''
\[
\widehat{A_{I}} \;\simeq\;
\underline{\mathrm{Hom}}^{!!}\Bigl(\prod_{j\in J^{-}} A_{I_{j}}^{\vee}\,,\
\prod_{j'\in J^{+}} \widehat{A_{I_{j'}}}^{\circ}\,;\ \mathrm{Ens}\Bigr)
\]
with $\approx$ arrows towards
\[
\underline{\mathrm{Hom}}_{!}\Bigl(\prod_{j\in J^{-}} A_{I_{j}}^{\vee},\
\widehat{A_{I^{+}}}\Bigr),
\qquad
\underline{\mathrm{Hom}}^{!}\Bigl(\prod_{j'\in J^{+}}
\widehat{A_{I_{j'}}}^{\circ},\ A_{I^{-}}^{\vee}\Bigr)
\]
\note{variances partly conjectural}

\section*{The duality between $\widehat{A}$ and $A^{\vee}$ (leaves 7 to 9)}

\leaf{7}
\note{author's page 3}
\[
\boxed{\ H_{A}(\Phi,\Psi)
= \mathrm{Hom}_{\widehat{A}}\bigl(\Phi,\ \Psi_{A}(\Psi)\bigr)
\simeq \mathrm{Hom}_{A^{\vee}}\bigl(\Psi,\ \varphi_{A}(\Phi)\bigr)\ }
\]
(in the margin: $\widehat{A}^{\circ} \xrightarrow{\varphi_{A}} A^{\vee}$,
$A^{\vee\circ} \xrightarrow{\Psi_{A}} \widehat{A}$,
$\widehat{A} \times A^{\vee\circ} \xrightarrow{H_{A}} \mathrm{Ens}$)

``i.e.\ one has $A^{\vee\circ} \rightleftarrows \widehat{A}$
($\varphi_{A}^{\circ}$, $\Psi_{A}$), a pair of adjoint functors.''

NB — a) the triangle
\begin{tikzcd}[column sep=small, row sep=small]
A^{\vee\circ} \arrow[rr, "\Psi_{A}"] & & \widehat{A} \\
& A \arrow[ul, hook] \arrow[ur, hook'] &
\end{tikzcd}
commutes, and $\Psi_{A}$ commutes with \uncertain{projective limits} — these
\uncertain{two} conditions \uncertain{characterise} $\Psi_{A}$ up to unique
isomorphism; — b) likewise for
$\widehat{A} \xrightarrow{\varphi_{A}^{\circ}} A^{\vee\circ}$:
``commutes, and $\varphi_{A}^{\circ}$ commutes with \ill{} and \ill{} — and
these conditions: $\varphi_{A}^{\circ}$ up to unique isomorphism.''

Left margin: \uncertain{$\varphi_{A^{\circ}} = \Psi_{A}$},
\uncertain{$\Psi_{A^{\circ}} = \varphi_{A}$},
\uncertain{$H_{A^{\circ}}(\Psi,\Phi) = H_{A}(\Phi,\Psi)$}
(\uncertain{formulas to be checked}).

Adjunction morphisms:
\[
\Phi \longrightarrow \Psi_{A}\,\varphi_{A}^{\circ}(\Phi)
\ \text{in } \widehat{A},
\qquad
\Psi \longrightarrow \varphi_{A}\,\Psi_{A}^{\circ}(\Psi)
\ \text{in } A^{\vee} ;
\]
``iso if $\Phi \in \mathrm{Ob}\,A$ \uncertain{or} $\Psi \in \mathrm{Ob}\,A$''.

c)
\[
\boxed{\ H_{A}(a,b) = \mathrm{Hom}_{A}(a,b)\ }
\]
``so $H_{A}$ extends $\mathrm{Hom}_{A}$, and is \uncertain{determined} up to
unique isomorphism by this property, together with the property that
$H_{A}$ commutes with \uncertain{limits} \ldots''
\[
H_{A}(\Phi, b) \simeq \mathrm{Hom}_{\widehat{A}}(\Phi,\ b),
\qquad
H_{A}(a, \Psi) \simeq \mathrm{Hom}_{A^{\vee}}(\Psi,\ a)
= \mathrm{Hom}_{A^{\vee\circ}}(a,\ \Psi)
\]
\note{second formula partly conjectural}

``Question: what to make of the adjunction morphisms $\alpha_{\Phi}$,
$\beta_{\Psi}$? Iso \uncertain{(yes)} only if $\Phi$ \uncertain{resp.}
$\Psi$ are representable? Is \uncertain{it} a weak equivalence (probably
not). Are the functors $\varphi_{A}$ and $\Psi_{A}$ \ill{}?''

``Relations with \uncertain{contractions}:''
\[
\boxed{\ H^{A} :\ (\Phi,\Psi) \longmapsto \Phi \star_{A} \Psi\ ;\quad
\widehat{A} \times A^{\vee} \longrightarrow \mathrm{Ens}\ }
\]
d)
\[
\boxed{\ H^{A}(a,b) = \mathrm{Hom}_{A}(b,a)\ }
\]
``and $H^{A}$ commutes with \uncertain{inductive limits}, under
\uncertain{which} conditions: up to unique isomorphism.''
\[
\boxed{\ \widehat{A} \xrightarrow{\ \sim\ }
\underline{\mathrm{Hom}}_{!}(A^{\vee}, \mathrm{Ens}),
\qquad
A^{\vee} \xrightarrow{\ \sim\ }
\underline{\mathrm{Hom}}_{!}(\widehat{A}, \mathrm{Ens})\ }
\]
\note{the decoration of the second Hom differs from the first — the shriek
sits differently, reading uncertain}
\[
\Phi \star a = \Phi(a), \qquad a \star \Psi = \Psi(a)
\]
\[
H^{A}(\Phi,\Psi) = \varphi^{A}(\Phi)(\Psi) \simeq \Psi^{A}(\Psi)(\Phi)
\]
\note{the first member is preceded by an unelucidated symbol}

\leaf{8}
\note{the whole leaf is struck out by three long diagonal strokes}
Multivariances:
\[
\Bigl(\prod_{i\in I} A_{i}\Bigr)^{\wedge},
\qquad
A_{1}^{\circ} \times A_{2}^{\circ} \times \cdots \times A_{n}^{\circ}
\to \mathrm{Ens},
\qquad
\widehat{A_{1}}^{\circ} \times \widehat{A_{2}} \times A_{3}^{\vee}
\times \cdots
\]
\note{the third list of variances is partly illegible}
\[
X_{1} \otimes X_{2} \simeq \mathrm{Hom}(X_{1}^{\vee}, X_{2})
\simeq \mathrm{Hom}(X_{2}^{\vee}, X_{1})
\]
\note{an intermediate member, carrying an underlined 1, remains uncertain}
\[
F(X_{1}\boxtimes X_{2})
\simeq \underline{\mathrm{Hom}}^{!!}\bigl(F(X_{1})^{\circ},\ F(X_{2})^{\circ};\
\mathrm{Ens}\bigr)
\simeq \cdots
\simeq \mathrm{Hom}\bigl(F(X_{1})^{\vee},\ F(X_{2})\bigr)
\simeq \mathrm{Hom}\bigl(F(X_{2})^{\vee},\ F(X_{1})\bigr)
\]
Families: $X_{I} = \prod_{j} X_{I_{j}}$, $I = I^{+} \amalg I^{-}$,
\[
F(X_{I})
\simeq \underline{\mathrm{Hom}}^{(1)}\Bigl(\prod_{j\in J} F(X_{I_{j}})^{\circ},\
\mathrm{Ens}\Bigr)
\simeq \underline{\mathrm{Hom}}_{(1)}\Bigl(\prod_{j\in J}
F(X_{I_{j}}^{\vee}),\ \mathrm{Ens}\Bigr)
\simeq \cdots
\]
\[
\cdots \simeq \underline{\mathrm{Hom}}_{!}\Bigl(\prod_{j\in J^{-}}
F(X_{I_{j}}),\ \prod_{j'\in J^{+}} F(X_{I_{j'}})\Bigr)
\simeq \underline{\mathrm{Hom}}^{!}\Bigl(\prod_{j'\in J^{+}}
F(X_{I_{j'}})^{\circ},\ \prod_{j\in J^{-}} F(X_{I_{j}}^{\vee})\Bigr)
\]
\note{decorations ① and exponents partly conjectural}
At the bottom: $\mathrm{Top}(A) \xrightarrow{\varphi} X \xleftarrow{\Psi}
\mathrm{Top}(B)$, \quad $A \to \mathrm{Pt}(X)$
\note{three final lines crossed out, mixing $\varphi_{*}(F)$,
$F \in \mathrm{Top}(A)$, $G \in \mathrm{Ob}\,B^{\vee}$ — reading too
uncertain to pursue}

\leaf{9}
\note{author's page 4; the leaf carries only these two boxed formulas}
\[
\bigl\langle \Phi_{*},\ f^{\circ*}(\Psi^{*}) \bigr\rangle_{A}
\simeq \bigl\langle f_{!}(\Phi_{*}),\ \Psi^{*} \bigr\rangle_{B},
\qquad
\bigl\langle f^{*}(\Psi_{*}),\ \Phi^{*} \bigr\rangle_{A}
\simeq \bigl\langle \Psi_{*},\ f^{\circ}_{!}(\Phi^{*}) \bigr\rangle_{B}
\]
\note{distribution of the stars and rings partly conjectural}

\section*{The envelope $\widetilde{A}$ (leaves 10 and 11)}

\leaf{10}
\note{author's page 5}
``$A \hookrightarrow B$ \uncertain{fully faithful}, $x \in B$:
\[
H_{x} \in \widehat{A}, \quad H_{x}(a) = \mathrm{Hom}_{B}(a, x) ;
\qquad
H^{x} \in A^{\vee}, \quad H^{x}(a) = \mathrm{Hom}_{B}(x, a)
\]
\[
H_{x}(a) \times H^{x}(b) \xrightarrow{\ \alpha_{x}\ } \mathrm{Hom}_{A}(a, b)
\]
''
\note{the arrow is composition in $B$; it lands in $\mathrm{Hom}_{A}$
because $A$ is full in $B$}

``Let $\widetilde{A}$ be the category of triples $(H_{*}, H^{*}, \alpha)$
where $H_{*} \in \widehat{A}$, $H^{*} \in A^{\vee}$,
$\alpha : H_{*}(a) \times H^{*}(b) \to \mathrm{Hom}_{A}(a,b)$ (functorial).
The arrows $(H_{*}, H^{*}, \alpha) \xrightarrow{\varphi}
(H'_{*}, H'^{*}, \alpha')$ being \uncertain{formed of the pairs}
$\varphi_{*} : H_{*} \to H'_{*}$ in $\widehat{A}$,
$\varphi^{*} : H'^{*} \to H^{*}$ in $A^{\vee}$, such that:
\[
\alpha_{a,b}\bigl(u,\ \varphi^{*}(u')\bigr)
= \alpha'_{a,b}\bigl(\varphi_{*}(u),\ u'\bigr)
\]
''
\note{the manuscript expresses this identity by a mixed-variance square,
$\varphi_{*}$ descending on the first factor and $\varphi^{*}$ ascending on
the second}

``For a \uncertain{fully faithful} functor $A \xrightarrow{\varphi} B$, one
\uncertain{deduces} a functor $\varphi^{\S} : B \to \widetilde{A}$,
$x \mapsto (\varphi^{\S}_{*}(x),\ \varphi^{\S *}(x),\ \alpha_{x})$.''
\note{decorations of the two components uncertain}

\leaf{11}
\note{author's page 6}
``NB $\widetilde{A}^{\circ} \simeq \widetilde{A}$ \ldots''
\note{end of line illegible}
``Case where $\varphi^{\S}$ is fully faithful? It suffices that $f$ be
\uncertain{strict.\ gén.} or \ill{} \uncertain{adjunctions}.''

``Case $A \xrightarrow{\mathrm{can}} \widehat{A}$:
$\varphi^{\S}(F)_{*} = F$,\quad
$\varphi^{\S}(F)^{*}(a) = \mathrm{Hom}(F, a)$;\quad $\varphi^{\S}$
\uncertain{faithful}.''

Small intermediate diagrams on $F$, $F^{\S}$, $G$, with $u \in F(a)$,
$u' \in F^{\S}(b)$.
\note{crossed-out sequence; one reads ``$v = v' \circ \varphi_{*}$!'' and
``we know that $v \circ u = v' \circ u'$''}

\struck{Let us then consider the full subcategory of $\widetilde{A}$ formed
by the} \ill{}
\note{passage struck out then rewritten, largely illegible; it contains the
chains $\Psi_{1}^{\circ} \to a \to \Phi \to \Psi_{2}^{\circ}$ and
$\Phi_{1} \to \Psi \to \Phi_{2}$}

$(F, G, \alpha) \simeq (F, F^{\S}, \alpha_{F})$, \qquad
$(F, G, \alpha) \simeq (G^{\S}, G, \alpha^{G})$
\note{conjectural reading of the two isomorphisms}

Functoriality: $A \xrightarrow{f} B$ induces
$\widetilde{f} : \widetilde{A} \to \widetilde{B}$; at the bottom of the
leaf, circled:
\begin{tikzcd}[column sep=normal, row sep=normal]
\widehat{A} \arrow[r, hook] \arrow[d, "f_{!}"']
  & \widetilde{A} \arrow[d, Rightarrow, "\widetilde{f}"]
  & A^{\vee\circ} \arrow[l, hook'] \arrow[d, "f^{\circ*\circ}"] \\
\widehat{B} \arrow[r, hook]
  & \widetilde{B}
  & B^{\vee\circ} \arrow[l, hook']
\end{tikzcd}
\note{labels of the vertical arrows partly conjectural; beside it, a column
of adjoint pairs — $f_{!}, f^{*}$; $f^{\circ}_{!}, f^{\circ *}$; \ldots{} —
too abbreviated to restore}

\section*{Self-dual core, gluings, coefficients (leaves 13 and 14)}

\leaf{13}
\note{author's page 7}
At the top:
\[
\widehat{A} \xrightarrow{\ \approx\ }
\underline{\mathrm{Hom}}_{!}(A^{\vee}, \mathrm{Ens}),
\qquad
A^{\vee\circ} \xrightarrow{\ \approx\ }
\underline{\mathrm{Hom}}_{\mathrm{rep}}(\widehat{A}, (\mathrm{Ens}))
\]
\note{the index of the second Hom — ``rep''? — is uncertain; above both, an
undecorated $\underline{\mathrm{Hom}}(A^{\vee}, (\mathrm{Ens}))$ receives
them}
At the centre, joined to $A$ by a vertical double arrow:
\[
\bigl(\mathrm{Kar}(A^{\circ})\bigr)^{\circ} \;\simeq\; \mathrm{Kar}(A)
\]
On the right, the triangle $\varepsilon(A)$ on $\widehat{A}$ and
$A^{\vee\circ}$, with
\[
A = \widehat{A} \cap A^{\vee\circ}
\]
\note{on the left of the leaf, an isolated $\overline{\varepsilon}(A)$,
below a separating stroke}

At the bottom, the warning about gluings:
``$N = M_{1} \cap M_{2}$ \uncertain{via} fully faithful arrows. Then one
\uncertain{defines} a category $\Sigma'$, \uncertain{union} of the full
subcategories $M'_{1}$ and $M'_{2}$, with \uncertain{intersection} $N'$,
and a fully faithful functor $\Sigma' \to M$ (\uncertain{factoring through}
$M_{1} \cup M_{2}$). But \uncertain{even} if $M_{1}$, $M_{2}$ are full in
$M$ (with $N$ \ill{} in $M$) and $M'_{1} \to M_{1}$, $M'_{2} \to M_{2}$
are equivalences, $M'_{1}$ and $M'_{2}$ are not \ill{} \ill{} full in
$\Sigma'$, and $N' \to N$ is not \ill{} an equivalence.''
\note{the reading of several words remains uncertain}

\leaf{14}
Duality with values in a category $M$:
\[
\underline{\mathrm{Hom}}(A, M)
\simeq \underline{\mathrm{Hom}}_{!}(\widehat{A}, M)
\]
\note{a third member, its decorations illegible, continues the line}
``Ex.\ $M = $ \ill{}: $\underline{\mathrm{Hom}}(B^{\circ}, M) \simeq \cdots$''
\note{example crossed out}

``$A = B^{\circ}$, $B = A^{\circ}$'':
\[
\underline{\mathrm{Hom}}(A, M) \simeq
\underline{\mathrm{Hom}}^{!}(\widehat{A}, M)
\]
``\uncertain{co}sheaves on $A$, ind.\ in $M$'';
\[
\underline{\mathrm{Hom}}(B^{\circ}, M) \simeq
\underline{\mathrm{Hom}}^{!}(\widehat{B}^{\circ}, M)
\]
``sheaves on $B$, ind.\ in $M$''.
\[
\mathcal{F}(X, M)^{\circ} \simeq \check{\mathcal{F}}(X, M^{\circ}),
\qquad
\bigl(\underline{\mathrm{Hom}}^{!}(A, M)\bigr)^{\circ}
\simeq \underline{\mathrm{Hom}}_{!}(A^{\circ}, M^{\circ})
\]
``$\mathcal{F}(X,M)$ and $\check{\mathcal{F}}(X, M')$ \uncertain{are deduced
from one another}'';
\[
\mathcal{F}(X, M) \approx \check{\mathcal{F}}(X^{\circ}, M),
\qquad
\underline{\mathrm{Hom}}^{!}(A^{\circ}, M)
\simeq \underline{\mathrm{Hom}}_{!}(A, M^{\circ})^{\circ}
\]
\note{final lines with uncertain decorations}

\end{document}
